%%%%%%%%%%%%%%%%%%%%%%%%%%%%%%%%%%%%%%%%%%%%%%%%%%%%%%%%%%
%%%%%%%%%%%%%%%%%%%%%%%%%%%%%%%%%%%%%%%%%%%%%%%%%%%%%%%%%%
%%%%%%%%%%%%%%%%%%%%%%%%%%%%%%%%%%%%%%%%%%%%%%%%%%%%%%%%%%
%%%%%%%%%%%%%%%%%%%%%%%%%%%%%%%%%%%%%%%%%%%%%%%%%%%%%%%%%%
\chapter{Installation and usage}

%%%%%%%%%%%%%%%%%%%%%%%%%%%%%%%%%%%%%%%%%%%%%%%%%%%%%%%%%%
%%%%%%%%%%%%%%%%%%%%%%%%%%%%%%%%%%%%%%%%%%%%%%%%%%%%%%%%%%
%%%%%%%%%%%%%%%%%%%%%%%%%%%%%%%%%%%%%%%%%%%%%%%%%%%%%%%%%%
\section{Requirements}

As of today, the library is tested against python \codeinline{3.8.9} on ubuntu \codeinline{20.04} and osx \codeinline{12.5}. The main required packages versions are the following:

\begin{minted}{python}
tensorflow==2.9.1
mpi4py==3.1.4
numpy==1.22.4
\end{minted}

The full list of required packages can be found in \codeinline{sparkle/requirements.txt}.

%%%%%%%%%%%%%%%%%%%%%%%%%%%%%%%%%%%%%%%%%%%%%%%%%%%%%%%%%%
%%%%%%%%%%%%%%%%%%%%%%%%%%%%%%%%%%%%%%%%%%%%%%%%%%%%%%%%%%
%%%%%%%%%%%%%%%%%%%%%%%%%%%%%%%%%%%%%%%%%%%%%%%%%%%%%%%%%%
\section{Installation}

Clone the repository locally, then install it using \codeinline{pip}:

\begin{minted}{python}
git clone git@github.com:jviquerat/sparkle.git
cd dragonfly
pip install -e .
\end{minted}

Once installed, the different tasks can be run directly from the command line through the \codeinline{spk} executable (see next section).

%%%%%%%%%%%%%%%%%%%%%%%%%%%%%%%%%%%%%%%%%%%%%%%%%%%%%%%%%%
%%%%%%%%%%%%%%%%%%%%%%%%%%%%%%%%%%%%%%%%%%%%%%%%%%%%%%%%%%
%%%%%%%%%%%%%%%%%%%%%%%%%%%%%%%%%%%%%%%%%%%%%%%%%%%%%%%%%%
\section{Usage}

%%%%%%%%%%%%%%%%%%%%%%%%%%%%%%%%%%%%%%%%%%%%%%%%%%%%%%%%%%
%%%%%%%%%%%%%%%%%%%%%%%%%%%%%%%%%%%%%%%%%%%%%%%%%%%%%%%%%%
\subsection{Training mode}

The training mode relies on the use of a \codeinline{.json} file for the description of the agent and its training mode (see section \textcolor{red}{to complete} for a complete description of the json parameter files):

\begin{minted}{python}
spk --train <filename>.json
\end{minted}

The completion of a training phase will produce a \codeinline{results/} repository in the current folder, in which training data and saved models will be stored (see section \textcolor{red}{to complete} for a description of the saved data). It will also produce a plot of the average and standard deviation of the score against training steps (see example in figure \ref{fig:training_score}).

%\input{fig/training_score}

%%%%%%%%%%%%%%%%%%%%%%%%%%%%%%%%%%%%%%%%%%%%%%%%%%%%%%%%%%
%%%%%%%%%%%%%%%%%%%%%%%%%%%%%%%%%%%%%%%%%%%%%%%%%%%%%%%%%%
\subsection{Exploitation mode}

Not available yet

%The exploitation mode requires the same \codeinline{.json} file used for training, as well as a model file. Model files are saved regularly during the training phase, and their saving/loading rely on the tensorflow \codeinline{save_weights} and \codeinline{load_weights} interfaces. The syntax to exploit a saved model is the following:
%
%\begin{minted}{python}
%dgf --eval -net <path/to/network/model> -json <filename>.json 
%\end{minted}
%
%In that case, the environment will just rely on the \codeinline{done} signal to stop the evaluation. Alternatively, you can provide a \codeinline{-steps <n_steps>} option, that will override the done signal of the environment, and force its execution for \codeinline{n_steps} steps:
%
%\begin{minted}{python}
%dgf --eval -net <path/to/network/model> -json <filename>.json -steps <n_steps>
%\end{minted}
%
%It is also possible to add a warmup phase, during which default control values are passed to the environment:
%
%\begin{minted}{python}
%dgf --eval -net ... -warmup <n_steps_warmup> <default_control_values>
%\end{minted}